\subsection{The Coordinate Formula for Euclidean Distance}

The Euclidean distance between two points is the length of the straight segment
joining them. Coordinates now allow this geometric length to be calculated.

For points
\[
P=(x_1,y_1),
\qquad
Q=(x_2,y_2),
\]
their projections show that the horizontal and vertical separations have
lengths
\[
|x_2-x_1|
\qquad\text{and}\qquad
|y_2-y_1|.
\]
Together with the segment $PQ$, these two perpendicular segments form a right
triangle. Here ``right'' refers to the Euclidean right angle already used in the
construction; no numerical measure of that angle is required.

\begin{center}
\begin{tikzpicture}[scale=0.9]
    \draw[axis] (-0.8,0) -- (5.2,0) node[right] {$x$};
    \draw[axis] (0,-0.8) -- (0,4.5) node[above] {$y$};
    \coordinate (P) at (1,1);
    \coordinate (Q) at (4,3.5);
    \coordinate (R) at (4,1);
    \draw[subset] (P) -- (Q);
    \draw[helper] (P) -- (R) node[midway,below] {$|x_2-x_1|$};
    \draw[helper] (R) -- (Q) node[midway,right] {$|y_2-y_1|$};
    \draw (3.72,1) -- (3.72,1.28) -- (4,1.28);
    \fill (P) circle (2pt) node[below left] {$P$};
    \fill (Q) circle (2pt) node[above right] {$Q$};
\end{tikzpicture}
\end{center}

By the Pythagorean theorem,
\[
d_E(P,Q)^2
=
(x_2-x_1)^2+(y_2-y_1)^2,
\]
and therefore
\[
\boxed{
d_E(P,Q)
=
\sqrt{(x_2-x_1)^2+(y_2-y_1)^2}
}.
\]

The formula does not create a new notion of distance. It is the numerical
coordinate expression of the Euclidean segment length already present in the
plane.
